%% threat_model.tex

\section{Threat Model and Motivation}
\label{sec:threat-model}

\subsection{System Model}

We consider a multi-agent LLM operating system (exemplified by \system{}) with the following architecture:
\begin{itemize}[leftmargin=*]
    \item \textbf{Agents} are scheduled as processes with per-agent \emph{roles} (e.g., \code{Code\_Reviewer}, \code{Python\_Coder}) drawn from a whitelisted blueprint registry.
    \item \textbf{Tools} (file I/O, shell, web fetch) are syscalls gated by a \code{PermissionChecker} that enforces role-based access control and tenant ownership.
    \item \textbf{Recovery orchestrator} classifies failures and dispatches a layered chain of recovery strategies organized into 11 conceptual tiers (from lightweight in-place retry up to circuit-breaker human escalation), realized by 14 concrete strategy classes---including reflection injection, topology mutation, model fallback, and human escalation. The 11/14 distinction is intentional: the tiers describe \emph{what kind} of failure is being handled (and at what escalation level), while the 14 classes are the concrete implementations, with some tiers containing multiple strategy variants (e.g., the resource-exhaustion tier includes both \code{ContextWindowRecovery} and \code{ResourceRefillStrategy}).
    \item \textbf{Semantic core dump} captures an agent's last input, context, and error trace upon crash, for post-mortem diagnosis by an LLM diagnostic module.
\end{itemize}

\subsection{Trust Boundaries}

The system maintains two trust boundaries (Figure~\ref{fig:trust-boundaries}):
\begin{enumerate}[leftmargin=*]
    \item \textbf{Entry-point boundary} (forward path): User input and external tool outputs are treated as untrusted. The \code{PermissionChecker}, \code{GuardrailEngine}, and \code{RecoveryPromptSanitizer} apply input validation at this boundary.
    \item \textbf{Recovery boundary} (reverse path): Error traces, core dumps, and diagnostic LLM outputs enter the next-round context \emph{without} re-validation. This boundary is the focus of this paper.
\end{enumerate}

\begin{figure}[htbp]
\centering
\begin{tikzpicture}[
    node distance=0.6cm and 0.7cm,
    scale=0.82, transform shape,
    box/.style={draw, rounded corners=2pt, minimum height=0.7cm, minimum width=2.0cm, align=center, font=\small},
    benign/.style={box, fill=green!15, draw=green!60!black},
    attack/.style={box, fill=red!15, draw=red!70!black},
    critical/.style={box, fill=red!70, draw=red!70!black, text=white, font=\small\bfseries},
    arr/.style={-Stealth, thick},
    barr/.style={-Stealth, thick, green!60!black},
    marr/.style={-Stealth, thick, red!70!black, dashed},
    region/.style={draw, dashed, rounded corners=4pt, inner sep=0.35cm}
]
% Benign flow
\node[benign] (user) {Legitimate User};
\node[benign, right=of user] (web1) {Normal Webpage};
\node[benign, right=of web1] (agent1) {Agent Executes};
\node[benign, right=of agent1] (done) {Task Done};
\draw[barr] (user) -- (web1);
\draw[barr] (web1) -- (agent1);
\draw[barr] (agent1) -- (done);

% Malicious flow
\node[attack, below=1.8cm of user] (atk) {Attacker\\(External Provider)};
\node[attack, right=of atk] (mal) {Malicious Webpage\\/ Poisoned API};
\node[attack, right=of mal] (fail) {Parsing Error\\(induced)};
\node[critical, right=of fail] (dump) {Poisoned Core Dump\\/ \code{lastErrorTrace}};
\draw[marr] (atk) -- (mal);
\draw[marr] (mal) -- (fail);
\draw[marr] (fail) -- (dump);

% Recovery path
\node[critical, right=of dump] (orch) {Recovery\\Orchestrator};
\node[critical, right=of orch] (esc) {Privilege\\Escalation};
\draw[marr] (dump) -- (orch);
\draw[marr] (orch) -- (esc);

% Region backgrounds
\begin{scope}[on background layer]
\node[region, fill=green!5, fit=(user)(web1)(agent1)(atk)(mal)(fail), label={[font=\footnotesize\itshape, green!60!black]above:Forward Path (entry-point defenses)}] {};
\node[region, fill=red!5, fit=(dump)(orch)(esc), label={[font=\footnotesize\itshape, red!70!black]above:Reverse Path (no re-validation)}] {};
\end{scope}
\end{tikzpicture}
\caption{Trust boundaries and attack chain for \rci{}. The forward path (top, green) applies entry-point validation; the reverse path (bottom, red) re-stamps externally-originated content with system-level trust, creating a trust-escalation trampoline.}
\label{fig:trust-boundaries}
\end{figure}

\subsection{Attacker Model}
\label{sec:attacker-model}

We assume the attacker is an \textbf{external content provider}---e.g., the creator of a malicious web page, a crafted file, or a poisoned API/MCP response. Concretely, the attacker controls \emph{external content} that the agent processes during normal operation:
\begin{itemize}[leftmargin=*]
    \item Web pages fetched via \code{web\_fetch} or \code{web\_scrape}
    \item Files read from the filesystem (if accessible)
    \item Tool outputs from external integrations (MCP servers, third-party APIs)
\end{itemize}

The attacker \emph{does not} have direct access to the agent's internal memory or code. In particular, the attacker cannot:
\begin{itemize}[leftmargin=*]
    \item Modify the agent's role, permission profile, or whitelisted blueprint
    \item Alter the LLM model, system prompt, or model provider
    \item Tamper with the recovery orchestrator's strategy dispatch logic or the \code{PermissionChecker}
    \item Inject code into the trusted computing base (kernel, runtime, or recovery machinery)
\end{itemize}

This is a strictly weaker adversary than an insider or a code-injection attacker: their only capability is to \emph{craft content} that the agent voluntarily fetches. The threat, therefore, must exploit a path by which such externally-fetched content \emph{acquires} system-level influence without ever touching the trusted computing base directly. The recovery channel provides exactly such a path.

\paragraph{Scope of the baseline threat model.}
The attacker model above---an external content provider who cannot influence tool registration or system configuration---is the \emph{baseline} threat studied in the main evaluation (\S\ref{sec:evaluation}). A stronger adversary who \emph{can} influence the tool-registration configuration (e.g., register an innocuously named MCP tool that fetches external content) is discussed as an \emph{extended} threat in \S\ref{sec:unregistered-tools}, along with the corresponding defense hardening (fail-closed default, explicit trust-class declaration, static-analysis guard). The main-evaluation ASR/FAR numbers are reported under the baseline threat model; the extended threat is addressed qualitatively and via the released code's fail-closed default, not re-measured end-to-end.

\subsection{Attack Chain}
\label{sec:attack-chain}

Recovery-channel injection follows a five-stage chain (Figure~\ref{fig:attack-chain}). The first three stages occur on the \emph{forward} path and are indistinguishable from benign operation; the last two occur on the \emph{reverse} path and constitute the actual injection.

\begin{enumerate}[label=(\roman*),leftmargin=*]
    \item \textbf{Content Delivery.} Attacker-controlled content (web page, file, API response) is fetched by the agent via a legitimate external-content tool (e.g., \code{web\_fetch}, \code{file\_read} on an untrusted path). At this stage the content is correctly treated as untrusted by the entry-point boundary.
    \item \textbf{Failure Induction.} The content is crafted to trigger a \emph{normal} tool parsing exception---malformed JSON, an unexpected HTTP 500, a syntax error in an embedded code snippet, or any failure mode that the recovery orchestrator routinely handles. No exotic bug is required; the orchestrator is designed to retry precisely these failures.
    \item \textbf{Payload Embedding.} The adversarial payload is embedded in the failure payload that the exception generates: the \code{lastErrorTrace} (for V1) or the \emph{semantic core dump} that captures the agent's last input and context on crash (for V2). Because the capture site does not record provenance, the payload is stored as an opaque \code{String}.
    \item \textbf{Recovery Trigger.} The orchestrator classifies the failure and dispatches a recovery strategy---\code{ReflectionInjectionRecovery} for V1, \code{TopologyMutationStrategy} for V2---that consumes the failure payload as input.
    \item \textbf{Privilege Escalation / Execution.} The payload is re-stamped with system-level trust and either (i)~injected into the next-round prompt under a \code{[SYSTEM CRITICAL]} template, causing the attacker's directive to be executed as a system command (V1), or (ii)~fed to a diagnostic LLM that emits an attacker-suggested high-privilege role, which reaches the node-replacement primitive without a permission re-check (V2).
\end{enumerate}

\begin{figure}[htbp]
\centering
\begin{tikzpicture}[
    node distance=0.6cm and 0.5cm,
    stage/.style={draw, rounded corners=2pt, minimum height=0.8cm, minimum width=2.2cm, align=center, font=\footnotesize},
    fwd/.style={stage, fill=green!12, draw=green!60!black},
    rev/.style={stage, fill=red!15, draw=red!70!black},
    final/.style={stage, fill=red!70, draw=red!70!black, text=white, font=\footnotesize\bfseries},
    arr/.style={-Stealth, thick, gray!70},
    lbl/.style={font=\tiny\itshape, gray!60!black, align=center}
]
% Stage chain
\node[fwd] (s1) {(i) Content\\Delivery};
\node[fwd, right=of s1] (s2) {(ii) Failure\\Induction};
\node[fwd, right=of s2] (s3) {(iii) Payload\\Embedding};
\node[rev, right=of s3] (s4) {(iv) Recovery\\Trigger};
\node[final, right=of s4] (s5) {(v) Privilege\\Escalation};

\draw[arr] (s1) -- (s2);
\draw[arr] (s2) -- (s3);
\draw[arr] (s3) -- (s4);
\draw[arr] (s4) -- (s5);

% Stage descriptions
\node[lbl, below=0.15cm of s1] {web\_fetch / file\_read};
\node[lbl, below=0.15cm of s2] {malformed JSON / HTTP 500};
\node[lbl, below=0.15cm of s3] {error trace / core dump};
\node[lbl, below=0.15cm of s4] {RecoveryOrchestrator};
\node[lbl, below=0.15cm of s5] {V1: prompt inject\\V2: role replace};

% Boundary marker
\draw[dashed, red!50, thick] ($(s3.east)+(0.25,-0.8)$) -- ($(s3.east)+(0.25,0.8)$);
\node[font=\tiny\itshape, red!50, right=0.05cm of s3.east, yshift=0.9cm] {trust boundary crossed};

% Region labels
\node[font=\footnotesize\itshape, green!60!black] at ($(s2.north)+(0,0.5)$) {Forward Path (benign-looking)};
\node[font=\footnotesize\itshape, red!70!black] at ($(s4.north)+(0.5,0.5)$) {Reverse Path (injection)};
\end{tikzpicture}
\caption{The five-stage \rci{} attack chain. Stages (i)--(iii) occur on the forward path and are indistinguishable from benign operation; stages (iv)--(v) occur on the reverse path and constitute the injection. The trust boundary is crossed between stages (iii) and (iv), where the recovery pipeline re-stamps the payload as system-trusted.}
\label{fig:attack-chain}
\end{figure}

The critical property of this chain is that \emph{no single stage is anomalous}: each stage is a feature of the self-healing design, and the attacker merely composes them. This is why forward-path defenses (input sanitization, tool-permission checks) cannot close the surface---the bypass occurs \emph{after} the forward path has legitimately accepted the content.

\subsection{Attack Vectors}

We identify two attack vectors via the recovery channel:

\textbf{Vector A: Reflection-Injection Trust Escalation (V1).} Adversarial content causes a tool failure; the error trace (containing the adversarial content) is injected into the next-round prompt under a high-trust \code{[SYSTEM CRITICAL]} template. The attacker's directive, now framed as a system command, is obeyed by the next-round agent.

\textbf{Vector B: Topology-Mutation Privilege Escalation (V2).} Adversarial content induces an agent crash; the core dump (containing the content) is fed to an LLM diagnostic module, which is tricked into emitting a high-privilege \code{suggested\_role}. The role replacement primitive executes without permission re-check, escalating privileges. Structurally, this is an instance of CWE-862 (Missing Authorization)~\cite{cwe862}: the role-replacement sink performs no authorization check on the role value it consumes, and the value is influenced by content that originated outside the trusted computing base. As discussed in \S\ref{sec:related-privilege}, the same configuration also realizes the classic confused-deputy pattern~\cite{hardy1988confused}: the recovery machinery (a privileged deputy) is induced by attacker-controlled content (a less-privileged caller) to exercise its role-replacement capability on the attacker's behalf. The two framings are complementary---CWE-862 identifies the missing check at the sink, while the confused-deputy framing identifies the authority-misuse pattern at the deputy---and both motivate the same structural fix (the \reauthgate{} privilege-downgrade invariant in \S\ref{sec:reauthgate}).

\subsection{The ``Trust-Escalation Trampoline''}

The core insight is that the recovery pipeline \emph{actively upgrades} the trust level of content it processes. Content that was correctly quarantined as untrusted on the forward path is re-laundered as trusted on the recovery path. This creates a \emph{trampoline}: the attacker does not need to defeat the entry-point defenses; they need only survive the recovery pipeline's (absent) re-validation.

Formally, let $c$ be content entering the forward path with trust label $T_{fwd}(c) = \textsc{Untrusted}$. Let $c'$ be the same content re-entering via the recovery path. The recovery pipeline assigns $T_{rec}(c') = \textsc{System}$ (by default, due to missing provenance tracking). The attacker exploits the gap: $T_{rec}(c') \gg T_{fwd}(c)$.
